% MedVision manuscript -- Related Work section

\section{Related Work}\label{sec:related-work}

\subsection{Medical image--text learning and report-derived targets}
Paired image--report resources support contrastive and local alignment methods
such as ConVIRT and GLoRIA, while clinical-domain language models improve the
handling of negation and uncertainty \citep{zhang2022convirt,huang2021gloria,boecking2022textsemantics}.
These methods establish the value of reports but do not show that a downstream
prediction depends on the image. MIMIC-CXR links radiographs to free-text
reports, whereas CheXpert derives uncertainty-aware labels from report language
\citep{johnson2019mimiccxr,irvin2019chexpert}. Report-labeling tools can also
propagate errors and subgroup differences \citep{sabottke2024targetselection,santomartino2024nlplabeling}.
The present audit therefore distinguishes report-conditioned prediction from
independent image-grounded diagnosis when both target and text input share a
report source.

\subsection{Shortcut learning and unimodal bias}
Shortcut-learning research shows that high in-distribution accuracy can arise from
predictive but unintended features \citep{geirhos2020shortcut}. In radiology,
acquisition devices, population structure, annotation practice, and clinical
context can create such shortcuts \citep{gichoya2023shortcuts}. In visual
question answering, VQA-CP, GVQA, product-of-experts methods, and RUBi reduce
answer or language priors by encouraging visual evidence \citep{agrawal2018dontassume,clark2019easyway,cadene2019rubi}.
They motivate LMH-style debiasing, but do not resolve the harder case in which
the benchmark label is itself text-derived.

\subsection{Auditing modality reliance}
Recent multimodal evaluations separate correctness from evidence dependence:
correct answers can be accompanied by flawed visual rationales, and stable
predictions need not be image-reliant \citep{jin2024hiddenflaws,sadanandan2026consistentdangerous}.
Chest-X-ray studies further report heterogeneity after controlling for clinical
context \citep{wang2026clinicalcontext}. Causal audits using image occlusion,
cross-patient replacement, and text-only comparisons provide a direct precedent
for testing image use \citep{lotfinia2026needimage}. Prototype and counterfactual
methods can improve saliency, but heatmaps alone cannot establish that image
removal changes a decision \citep{chiang2026cxrnexus}.

\subsection{Fusion gates and explanation limits}
Gated multimodal units learn modality weights, yet a non-zero or balanced gate
is not causal attribution \citep{arevalo2017gmu}. A gate may remain near its
initialization when the objective does not identify a different value. Likewise,
reasoning traces or plausible radiologist-style explanations are useful
interfaces but not grounding tests \citep{myronenko2025reasoningcxr}. The
present study combines checkpoint gate extraction with a report-preserving
image intervention and treats Grad-CAM as a qualitative sanity check.

\subsection{Positioning of the present study}
Prior work treats image--text learning, report-label uncertainty, shortcut bias,
and behavioural audits separately. We examine their intersection using
report-preserving, text-only, and image-only controls, an operational ICM,
checkpoint gates, and seed-level inference. The study tests incremental visual
evidence without claiming that report-conditioned prediction or all medical
VLMs are invalid.
